\documentclass[10pt,a4paper]{amsart}
\usepackage[margin=19mm]{geometry}
\usepackage{amsmath,amssymb,mathtools}
\usepackage[scr=rsfs,cal=euler]{mathalfa}
\usepackage{xfrac}
\usepackage[unicode,hidelinks,bookmarks=false]{hyperref}
\newcommand{\ie}{i.e.\ }
\newcommand{\cf}{cf.\ }
\newcommand{\resp}{respectively\ }
\newcommand{\Subsection}{\subsection}
\providecommand{\nobreakdash}{\nobreak\hbox{-}}
\setlength{\emergencystretch}{2em}
\multlinegap0pt
\usepackage{bm}
\usepackage{bbm}
\usepackage{dsfont}
\usepackage{xspace}
\usepackage{cancel}
\usepackage{mathtools}
\usepackage{amsthm}
\makeatletter
\@ifundefined{thm}{\newtheorem{thm}{Thm}}{}
\@ifundefined{theorem}{\newtheorem{theorem}[thm]{Theorem}}{}
\makeatother
\newcommand\E{\operatorname{\mathbf{E}}}
\title[Mean-Field Games with two-sided singular controls for L\'evy]{Mean-Field Games with two-sided singular controls for L\'evy processes}
\author{Facundo Oli\'u}
\date{Source: arXiv:2505.22936v1 (CC BY 4.0)}
\begin{document}
\maketitle

\noindent\emph{Source and licence.} Mean-Field Games with two-sided singular controls for L\'evy processes. Version arXiv:2505.22936v1 (CC BY 4.0), distributed under the Creative Commons Attribution 4.0 International licence, as recorded in the arXiv licence metadata for this exact version and shown on the abstract page https://arxiv.org/abs/2505.22936v1. Full rights notice: licenses/arxiv-2505.22936-CC-BY-4.0.txt.\par\smallskip

\noindent\textit{Editorial scope.} Counted unit: the Theorem whose source label is \texttt{None} in the article (source0.tex lines 1033--1042, source label \texttt{None}), delivered together with the article's own complete original proof of it (source0.tex lines 1044--1089, one \texttt{proof} environment of 46 lines). No mathematics is written, repaired or re-derived: statement and proof are the article's own text line for line. Required context retained uncounted: none. Every definition, display and label of those lines is retained unchanged. Omitted from this package and not redistributed: 1--1032, 1043--1043, 1090--1465. Nothing inside a retained excerpt is omitted or altered: every retained body is delivered exactly as the article's lines are written, so no omission receipt of any kind accompanies this package. Citations: the retained excerpts carry no citation command at all, so no bibliography entry of the article is transcribed and no reference is redistributed. Reference adaptation: none was needed and none was made; every reference inside the delivered unit resolves inside it, so no unresolved reference or citation is produced. Printed slips kept literal: no printed word, symbol, hypothesis or number of the counted statement or of its proof was altered. Layout and class: the article's own class is \texttt{article} with packages including natbib, float, xcolor, graphicx, hyperref, amsmath, amsthm, amssymb, amsfonts, bigints, wrapfig, caption, subcaption, mathtools; the wrapper is the lane's pinned \texttt{amsart} editorial wrapper at 10pt on A4, which restates the theorem-like environments the retained text needs and adds the article's own macro definitions that the retained text needs (\texttt{\textbackslash E}), with extra wrapper packages (bm, bbm, dsfont, xspace, cancel, mathtools). The delivered numbering is the unit's own. Assets: no retained span contains a figure environment, a graphics-inclusion command, a PDF-page inclusion, a table or a TikZ picture; no image file is copied into this package, the package declares no asset dependency and no image file is redistributed with it. Licence: version arXiv:2505.22936v1 of this article is distributed under the Creative Commons Attribution 4.0 International licence, as recorded in the arXiv licence metadata for this exact version and shown on the abstract page https://arxiv.org/abs/2505.22936v1. A full rights notice is delivered at licenses/arxiv-2505.22936-CC-BY-4.0.txt. Characters: every retained excerpt is pure ASCII, so the delivered engine, XeLaTeX with its default Latin Modern text font, reports no missing character.
\begin{theorem}
    Fix $N \in \mathbb{N}, \ K>0,  $ and $a < b$. Let the constant $\delta>0$ satisfy $\vert f(x)-f(y) \vert <1/N^2 $.
 \begin{itemize}
     \item[\rm{(i)}] Assume that the set of admissible controls is restricted to $\mathcal{B}^K_0$. If $(a,b)$ is an ergodic MFG-equilibrium and $\Lambda_{a,b} \in \mathcal{B}^K_0$, then $\Lambda_{a,b}$ is an $r$-ergodic Nash equilibrium for 
     $$r=\max_{z \in [a,b]}(h^2(z)) 4K e^{\frac{-2 \delta^2}{(b-a)^2N}}+2K/N .$$
     \item[\rm{(ii)}] Assume that the set of admissible controls is restricted to $\mathcal{B}^K_\epsilon$. If $(a,b)$ is an $\epsilon$-discounted MFG equilibrium and $\Lambda_{a,b} \in \mathcal{B}^K_{\epsilon}$, then $\Lambda_{a,b}$ is an $r,\epsilon$-discounted Nash equilibrium for
          $$r=\frac{\max_{z \in [a,b]}(h(z))^2 4Ke^{\frac{-2 \delta^2}{(b-a)^2N}}}{\epsilon}+\frac{2K}{N\epsilon} . $$

 \end{itemize}   
\end{theorem}

\begin{proof}
    We only prove the first claim as the other is analogue.
    First observe, as $f(x)$ is continuous, the set $f([a,b])$ is a closed interval, denote it by $[m,M]$,
and  observe that 
$$
(X^{i,a,b}_s,\bar f^{-i,a,b}_s)\in [a,b]\times[m,M], 
$$
that is a product of closed intervals. 
As for $m \leq y \leq M$, we have   
%
\begin{equation}\label{eq:cbounded}
    c(x,y) \leq g(x)\max_{z \in [m,M]}h(z)  =: g(x)H(m,M)  ,
\end{equation}
we define
 $$C:= \lim_{T \to \infty} 2 \left( \frac{H(m,M)}{T}\E_x \left( \int_0^T g ( X^{a,b}_s ) ds +dU_T^{a,b}q_u+dD_T^{a,b}q_d \right) \right) . $$
%
  Let $\mathcal{C}$ be defined as the set of strategies $\eta \in \mathcal{B}^K_0$ that satisfy 
  $$J_{\infty,N}^i (\eta, \Lambda^{-i}_{a,b})<C , \quad \text{ for all } i=1, \dots N, \ N \in \mathbb{N}. $$
 Now, observe that for all $i=1, \dots N, \ N \in \mathbb{N}$, by taking $\eta = (U^{a,b},D^{a,b})$, we have  $ J_{\infty,N}^i (\eta,\Lambda^{-i})\leq C/2< C ,$ thus
\begin{equation*}
    \inf_{\eta \in \mathcal{B}^K_0} J_{\infty,N}^i (\eta,\Lambda^{-i}_{a,b})=   \inf_{\eta \in \mathcal{C}} J_{\infty,N}^i (\eta,\Lambda^{-i}_{a,b}), \quad \text{ for all } i=1, \dots N, \ N \in \mathbb{N}.  
\end{equation*}
Furthermore, observe that

\begin{align}
J_{\infty,N}^{i}( \Lambda_{a,b})&=  J_{\infty,N}^{i}(\eta,\Lambda^{-i}_{a,b})+\big(J_{\infty,N}^{i}(\Lambda_{a,b})-J(\eta,X^{a,b}_{\infty}\big) \nonumber
\\& +\big(J(\eta,X^{a,b}_{\infty}\big)-J_{\infty,N}^{i}(\eta,\Lambda^{-i}_{a,b}) \big)\nonumber \\
                &\leq J_{\infty,N}^i(\eta,\Lambda^{-i}_{a,b})+\big(J_{\infty,N}^{i}(\Lambda_{a,b})-J((U^{a,b},D^{a,b}),X^{a,b}_{\infty}\big)\nonumber\\&+\big(J(\eta,X^{a,b}_{\infty} )-J_{\infty,N}^{i}(\eta,\Lambda^{-i}_{a,b}) \big).  \nonumber
\end{align}
%
Thus, to prove (\rm{i}), it is enough to demostrate 
 \begin{equation}\label{E:NaproximmationSubstraction2}
\sup_{\eta\in\mathcal{C}}\vert  J_{\infty,N}^{i}(\eta, \Lambda^{-i}_{a,b})-J(\eta, X^{a,b}_{\infty})\vert  \leq r/2,
\end{equation}
as the term $\vert J_{\infty,N}^{i} (\Lambda_{a,b})- J((U^{a,b},D^{a,b}),X^{a,b}_{\infty}) \vert$ is smaller or equal than the supremum defined in \eqref{E:NaproximmationSubstraction2}, because $(U^{a,b},D^{a,b}) \in \mathcal{C}$. For that endeavor, we need to prove that
\begin{equation}\label{eq:finalapproximation}
\left|\frac1T\E_x\int_0^T\left(c(X^{i,\eta_i}_s,\bar f^{a,b,-i})-c(X^{i,\eta_i}_s,\E_x(f(X^{a,b}_{\infty}))\right)\,ds\right|,
\end{equation}
is smaller or equal to $r/2$.
With Cauchy-Schwarz inequality and the fact that $\eta^i \in \mathcal{B}_0^K$, we deduce the expression \eqref{eq:finalapproximation} is smaller or equal than 
\begin{equation}\label{eq:finalapproximation2}
\left|\frac{K}{T}\E_x\int_0^T\left(h(\bar f^{a,b,-i})-h(\E_x(f(X^{a,b}_{\infty}))\right)^2\,ds\right|.
\end{equation}
Finally, we use Hoeffding's inequality for bounded random variables, the fact $\vert h(x)-h(y)\vert \leq b^2+a^2$ and the definition of $\delta$ to conclude \eqref{eq:finalapproximation2} is smaller or equal than $r/2$, thus concluding the proof of the theorem.

\end{proof}

\end{document}
